\chapter{Instrukcje montażu}\label{ch:montaz}

Montuj zestaw według tej kolejności. Każdy podzespół ma osobę odpowiedzialną (formularz F1). Po każdym kroku zaznacz go w formularzu.

\begin{sekcja}{Wkładka gondoli}
Wkładka to płyta montażowa na dnie pudełka. Do wkładki mocujesz całą elektronikę.
\begin{enumerate}
  \item Wytnij płytę z lekkiego i wytrzymałego materiału: plexi, sklejka albo włókno węglowe. Wymiar: wewnętrzne dno pudełka.
  \item Rozplanuj elementy tak, aby środek masy był na środku płyty.
  \item Wywierć otwory na śruby i opaski zaciskowe. Przyklej rzepy pod lekkie elementy.
  \item Zamocuj płytę w pudełku (klej do styropianu albo śruby z podkładkami przez dno).
\end{enumerate}

\begin{uwaga}Uderzenie przy lądowaniu jest bardzo silne. Każdy element musi być przymocowany do wkładki. Karta pamięci musi być zabezpieczona przed wypadnięciem (taśma na gnieździe). Utrata danych to utrata całej misji.\end{uwaga}
\end{sekcja}

\begin{sekcja}{Gondola platformy}
\begin{enumerate}
  \item Zważ puste pudełko.
  \item Zamocuj do wkładki: HABduino, OBC, baterie, SPOT, znaczniki Bluetooth.
  \item Poprowadź anteny nadajników przez ścianę albo dno pudełka. Anteny skierowane w dół.
  \item Zamocuj antenę GPS pod pokrywą albo na pokrywie, z dala od baterii i metalu.
  \item Zamocuj płytę montażową na pokrywie. Zamocuj zwalniacz serwomechaniczny. Poprowadź przewód serwa do środka.
  \item Podłącz baterie. Sprawdź diody statusu.
  \item Sprawdź odbiór telemetrii z odległości co najmniej \SI{100}{m}.
  \item Sprawdź zwolnienie odcinacza na komendę.
  \item Sprawdź równowagę: zawieś gondolę na linkach. Gondola wisi poziomo.
  \item Wyłącz zasilanie. Zważ gondolę. Wpisz masę do budżetu masy.
\end{enumerate}
\end{sekcja}

\begin{sekcja}{Olinowanie}
\begin{enumerate}
  \item Przytnij taśmy na długość. Zamontuj regulatory.
  \item Połącz gondolę ładunku z gondolą platformy taśmą (dookoła obu pudełek, na krzyż).
  \item Przygotuj krótką pętlę do szyjki powłoki z krętlikiem.
  \item Przywiąż spadochron pod krętlikiem.
  \item Przygotuj zapas linki 10--15~m między spadochronem a gondolą. Zwiń zapas na kartonie albo rurce, aby się nie splątał.
  \item Przygotuj linkę odcinacza: od krętlika do zwalniacza w gondoli.
  \item Zważ olinowanie ze spadochronem. Wpisz masę do budżetu masy.
\end{enumerate}
\end{sekcja}

\begin{sekcja}{Kontrola w warsztacie}
\begin{itemize}
  \item[\kratka] Łączna masa ładunku poniżej \SI{4}{kg}.
  \item[\kratka] Każdy pakiet poniżej \SI{2}{kg}.
  \item[\kratka] Opis markerem na każdym pakiecie: telefon, „Projekt badawczy”, logo.
  \item[\kratka] Wszystkie testy z formularza F5 zaliczone.
  \item[\kratka] Wszystko spakowane według wykazu materiałów.
\end{itemize}
\end{sekcja}

\begin{sekcja}{Na stanowisku startowym}
\begin{enumerate}
  \item Rozłóż podłoże: plandeka, koc, folia dolna. Dociśnij rogi skrzynkami.
  \item Rozłóż zestaw lotny w linii, z wiatrem: krętlik, spadochron, zapas linki, gondole.
  \item Włącz elektronikę platformy i ładunku. Sprawdź pozycję GPS i odbiór telemetrii na mapach.
  \item Zatankuj powłokę (rozdział~\ref{ch:naziemny}).
  \item Zamknij szyjkę. Przywiąż pętlę z krętlikiem.
  \item Zaczep linkę odcinacza o zwalniacz. Uzbrój odcinacz.
  \item Sprawdź węzły. Zabezpiecz je taśmą.
  \item Zamknij gondole. Owiń je taśmą. Napnij pasy.
  \item Wykonaj formularz F11 „Start”.
\end{enumerate}
\end{sekcja}

\begin{sekcja}{Start}
\begin{enumerate}
  \item Pomocnicy trzymają powłokę folią górną. Jedna osoba trzyma gondolę.
  \item Kierownik startu sprawdza niebo (samoloty) i odlicza.
  \item Na komendę pomocnicy zdejmują folię i puszczają powłokę. Osoba z gondolą idzie pod wiatr, aż linki się naprężą.
  \item Osoba z gondolą puszcza gondolę, gdy gondola jest pod balonem. Nie rzucaj gondoli.
  \item Zapisz czas startu.
\end{enumerate}

\dwafoto{2020_start_torun}{Start: powłoka unosi linkę, zespół trzyma gondolę (20.07.2020).}
{2023_zestaw_przed_startem}{Zestaw rozłożony z wiatrem przed startem (26.05.2023).}
\end{sekcja}
